% ================================================================
%  PREAMBULO – Universidad de los Andes  |  Taller 2
%  Basado en la plantilla de documentos (Parcial 1, v8.0)
%  Compilar desde la carpeta del documento: latexmk -pdf main.tex
% ================================================================

% ── IDIOMA ───────────────────────────────────────────────────
\usepackage[utf8]{inputenc}
\usepackage[T1]{fontenc}
\usepackage{newtxtext, newtxmath}   % Times: trazo más grueso y legible
\usepackage[spanish, es-noshorthands]{babel}
\addto\captionsspanish{
  \renewcommand{\tablename}{Tabla}
  \renewcommand{\listtablename}{Índice de Tablas}
  \renewcommand{\listfigurename}{Índice de Figuras}
  \renewcommand{\contentsname}{Tabla de Contenido}
}

% ── MÁRGENES ─────────────────────────────────────────────────
\usepackage[
  top=2.5cm, bottom=2.5cm,
  left=2.5cm, right=2.5cm,
  headheight=14pt
]{geometry}

% ── ENCABEZADO Y PIE ─────────────────────────────────────────
\usepackage{fancyhdr}
\usepackage{lastpage}

\newcommand{\estilopagina}{%
  \fancyhf{}%
  \renewcommand{\headrulewidth}{0.4pt}%
  \renewcommand{\footrulewidth}{0.4pt}%
  \fancyhead[L]{\small\textit{\DocCurso}}%
  \fancyhead[R]{\small\textit{\DocInstitucion}}%
  \fancyfoot[L]{\small\DocEntrega}%
  \fancyfoot[R]{\small\textbf{Pág.~\thepage~/~\pageref{LastPage}}}%
}
\pagestyle{fancy}
\estilopagina
% Los índices usan el estilo plain: mismo encabezado y pie
\fancypagestyle{plain}{\estilopagina}

% ── PÁGINAS HORIZONTALES ─────────────────────────────────────
\usepackage{pdflscape}
\usepackage{tikz}
\usepackage{everypage}

\newif\ifinlscape
\inlscapefalse

% Encabezado y pie para páginas horizontales (arriba y abajo en vista landscape)
\AddEverypageHook{%
  \ifinlscape
    \begin{tikzpicture}[remember picture, overlay]
      % Encabezado en la parte superior de la página horizontal
      \node[rotate=90, anchor=center, inner sep=2pt]
        at ([xshift=1.5cm]current page.west)
        {\parbox{\dimexpr\paperheight-5cm\relax}%
         {\small\itshape\DocCurso\hfill\DocInstitucion}};
      % Pie en la parte inferior de la página horizontal
      \node[rotate=90, anchor=center, inner sep=2pt]
        at ([xshift=-1.5cm]current page.east)
        {\parbox{\dimexpr\paperheight-5cm\relax}%
         {\small\DocEntrega\hfill\textbf{Pág.~\thepage~/~\pageref{LastPage}}}};
    \end{tikzpicture}%
  \fi%
}

\newenvironment{pagehorizontal}{%
  \begin{landscape}%
  \inlscapetrue%
  \pagestyle{empty}%
}{%
  \end{landscape}%
  \inlscapefalse%
  \pagestyle{fancy}%
}

% ── MATEMÁTICAS ──────────────────────────────────────────────
\usepackage{amsmath}

% ── GRÁFICOS Y TIKZ ──────────────────────────────────────────
\usepackage{graphicx}
\usepackage{float}
\graphicspath{{../assets/}{./}}
\usetikzlibrary{shapes, shapes.multipart, shapes.misc, arrows, arrows.meta,
                positioning, shadows, calc, fit}

% ── TABLAS ───────────────────────────────────────────────────
\usepackage{booktabs}
\usepackage{tabularx}
\usepackage{multirow}
\usepackage{multicol}
\usepackage{longtable}
\usepackage{array}
\usepackage{colortbl}
\usepackage{adjustbox}

\newcolumntype{L}[1]{>{\raggedright\arraybackslash}p{#1}}
\newcolumntype{C}[1]{>{\centering\arraybackslash}p{#1}}
\newcolumntype{R}[1]{>{\raggedleft\arraybackslash}p{#1}}

\renewcommand{\arraystretch}{1.2}
\setlength{\tabcolsep}{4pt}

% Encabezado de fila oscuro
\newcommand{\tblhead}{\rowcolor[gray]{0.25}\color{white}\bfseries}

% ── CAPTIONS ─────────────────────────────────────────────────
\usepackage{caption}
\captionsetup{font=small, labelfont=bf, labelsep=period, skip=3pt}
\usepackage{chngcntr}
\counterwithout{table}{section}
\counterwithout{figure}{section}

% ── SECCIONES ────────────────────────────────────────────────
\usepackage{titlesec}
\titleformat{\section}
  {\normalsize\bfseries}{\thesection.}{0.5em}{}
  [\vspace{1pt}\rule{\linewidth}{0.4pt}]
\titleformat{\subsection}
  {\normalsize\bfseries}{\thesubsection.}{0.5em}{}
\titleformat{\subsubsection}
  {\normalsize\itshape}{\thesubsubsection.}{0.4em}{}
\titlespacing{\section}      {0pt}{10pt}{4pt}
\titlespacing{\subsection}   {0pt}{8pt}{3pt}
\titlespacing{\subsubsection}{0pt}{6pt}{2pt}

% ── ÍNDICE ───────────────────────────────────────────────────
\usepackage{tocloft}
\setlength{\cftbeforesecskip}{3pt}
\setlength{\cftbeforesubsecskip}{1pt}

% ── LISTAS ───────────────────────────────────────────────────
\usepackage{enumitem}
\setlist[itemize]  {leftmargin=1.2em, itemsep=0pt, topsep=2pt, parsep=0pt}
\setlist[enumerate]{leftmargin=1.2em, itemsep=0pt, topsep=2pt, parsep=0pt}

% ── CAJAS ────────────────────────────────────────────────────
\usepackage[many]{tcolorbox}
\newtcolorbox{notebox}[1][]{
  enhanced, breakable, arc=0pt,
  colback=white, colframe=black,
  leftrule=2pt, toprule=0pt, bottomrule=0pt, rightrule=0pt,
  fonttitle=\bfseries\small, title={#1}, coltitle=black, colbacktitle=white,
  left=6pt, right=4pt, top=3pt, bottom=3pt,
  before skip=4pt, after skip=4pt
}
% Contenido por redactar; buscar \pendiente antes de entregar
\newcommand{\pendiente}[1]{\begin{notebox}[Pendiente]\small #1\end{notebox}}

% ── CÓDIGO Y CONTRATOS ───────────────────────────────────────
\usepackage{listings}
\renewcommand{\lstlistingname}{Listado}
\lstdefinelanguage{json}{morestring=[b]", sensitive=true}
\lstset{
  basicstyle=\ttfamily\footnotesize,
  columns=fullflexible, keepspaces=true,
  breaklines=true, showstringspaces=false,
  frame=leftline, xleftmargin=8pt, framexleftmargin=6pt,
  captionpos=b, upquote=true,
  extendedchars=true,
  literate={á}{{\'a}}1 {é}{{\'e}}1 {í}{{\'i}}1 {ó}{{\'o}}1 {ú}{{\'u}}1
           {Á}{{\'A}}1 {É}{{\'E}}1 {Í}{{\'I}}1 {Ó}{{\'O}}1 {Ú}{{\'U}}1
           {ñ}{{\~n}}1 {Ñ}{{\~N}}1 {ü}{{\"u}}1 {¿}{{?`}}1 {¡}{{!`}}1
}

% ── ADR ──────────────────────────────────────────────────────
% \begin{adr}{ADR-01}{Título}{Aceptada} \adrcampo{Contexto} texto ... \end{adr}
\newenvironment{adr}[3]{%
  \subsection{#1: #2}%
  \noindent\textbf{Estado:} #3\par
}{\par\medskip}
\newcommand{\adrcampo}[1]{\par\smallskip\noindent\textbf{#1.}\ }

% ── HIPERVÍNCULOS ────────────────────────────────────────────
\usepackage{hyperref}
\hypersetup{
  hidelinks,
  pdftitle  = {\DocTitulo},
  pdfauthor = {\DocGrupo},
  pdfsubject= {\DocCurso}
}

% ── OTROS ────────────────────────────────────────────────────
\usepackage{csquotes}
\newcommand{\secdiv}{\vspace{3pt}\rule{\linewidth}{0.3pt}\vspace{3pt}}
